%! AP Statistics — Unit 4, Lesson 4.1 — Student Guided Notes
\documentclass[10pt]{article}
\usepackage{apstats-article}
\usepackage{apstats-boxes}

\begin{document}

\pageheader{Unit 4, Lesson 4.1}{Guided Notes}
\namedateperiod

\vspace{0.12in}

\begin{objectivebox}
\small By the end of this lesson, I will be able to\dots\\[4pt]
\writeline\\[1pt]
\writeline\\[1pt]
\writeline
\end{objectivebox}

\vspace{0.1in}

\begin{vocabbox}
\termblanklong{Random process}
\termblanklong{Pattern}
\termblanklong{Streak}
\termblanklong{Random variation}
\termblanklong{Statistical question}
\end{vocabbox}

\vspace{0.1in}

\begin{hookbox}
\small\textbf{Coin flip sequences:}
\begin{itemize}[leftmargin=1.5em, itemsep=3pt]
  \item Sequence A: H H T H T T H H H T T H H T H T A T H T H
  \item Sequence B: H H H H H T T T T T H H H H H T T T T T
  \item Sequence C: H T H T T H T H H H T T H T H T H H T H
\end{itemize}
\textbf{Which looks most like a real fair coin? Which looks suspicious? Why?}\\[4pt]
\writeline\\[1pt]\writeline
\end{hookbox}

\vspace{0.1in}

\begin{notesbox}{Patterns in Random Processes}
\small
\textbf{Key Principle (VAR-1.F.1):} Patterns in data \blank{4cm} necessarily mean
that the variation is \blank{4cm}.

\vspace{10pt}
\renewcommand{\arraystretch}{1.6}
\begin{tabularx}{\linewidth}{@{} >{\bfseries\color{navy}}p{0.3\linewidth} X @{}}
Random process & \blank{8cm} \\
Pattern & \blank{8cm} \\
Streak & \blank{8cm} \\
\end{tabularx}

\vspace{10pt}
\textbf{The Hot-Hand Example:}\\
A basketball player makes 7 consecutive free throws (career average: 70\%).
\begin{itemize}[itemsep=4pt]
  \item The pattern suggests: \blank{7cm}
  \item But we should ask: \blank{7cm}
  \item To answer that question, we need: \blank{5cm}
\end{itemize}

\vspace{8pt}
\textbf{Counter-intuitive fact:} Most people expect random sequences to have
\blank{4cm} alternations. In reality, a run of 5+ identical outcomes in
20 coin flips has about a \blank{2cm} chance of occurring.
\end{notesbox}

\vspace{0.1in}

\begin{notesbox}{Identifying Statistical Questions from Patterns}
\small
A \textbf{statistical question} is one that \blank{6cm}.

\vspace{8pt}
For any pattern in data, ask these questions:
\vspace{4pt}
\renewcommand{\arraystretch}{1.8}
\begin{tabularx}{\linewidth}{@{} l X @{}}
\textbf{Question} & \textbf{My answer} \\
\hline
What pattern do I observe? & \blank{8cm} \\
What claim is the pattern tempting me to make? & \blank{8cm} \\
Could chance alone produce this pattern? & \blank{8cm} \\
What evidence would convince me it is NOT random? & \blank{8cm} \\
\end{tabularx}
\end{notesbox}

\vspace{0.1in}

\begin{practicebox}
\small For each scenario, (i) describe the pattern and (ii) write a statistical question it suggests.

\vspace{6pt}
\begin{enumerate}[leftmargin=1.5em, itemsep=12pt]
  \item A small town reports 8 cases of a rare disease in one year when the expected rate
        predicts fewer than 2.
        \begin{enumerate}[label=(\alph*), itemsep=4pt]
          \item Pattern: \blank{9cm}
          \item Statistical question: \writeline
        \end{enumerate}

  \item A coffee shop notices that customers order hot drinks on 9 of the last 10 rainy days.
        \begin{enumerate}[label=(\alph*), itemsep=4pt]
          \item Pattern: \blank{9cm}
          \item Statistical question: \writeline
        \end{enumerate}

  \item A student flips a coin 15 times and gets heads 11 times.
        \begin{enumerate}[label=(\alph*), itemsep=4pt]
          \item Pattern: \blank{9cm}
          \item Could this happen by chance with a fair coin? \blank{4cm}
        \end{enumerate}
\end{enumerate}
\end{practicebox}

\end{document}
